\chapter{Algebraic Curves}
\begin{center}\itshape
A line meets a parabola twice, a circle meets a circle twice, a cubic meets a cubic nine times.\\
Except when they do not.\\
Count at infinity, count over $\C$, count with multiplicity, and the exceptions disappear.
\end{center}

In 1720 Colin Maclaurin noticed something odd. A cubic curve in the plane has an equation with ten coefficients, so nine points should
determine it, just as two points determine a line and five a conic. But two different cubics meet in nine points. So nine points do
\emph{not} always determine a cubic, and Gabriel Cramer, who made the same observation in 1750, could not say which nine points do.
The ``Cramer paradox'' sat there until the nineteenth century, when it was understood as the first instance of two theorems.
\emph{B\'ezout's theorem} (stated by Newton and Maclaurin, ``proved'' by B\'ezout in 1779, proved by Halphen and Noether a century later):
two plane curves of degrees $m$ and $n$ without a common component meet in exactly $mn$ points, if one counts properly. And the
\emph{Cayley--Bacharach theorem}: the nine points where two cubics meet are special, since any cubic through eight of them passes through
the ninth.

Counting properly needs three things. Points at infinity: the parallel lines $y=0$ and $y=1$ meet there. Complex points: the circle
$x^2+y^2=1$ and the line $x=2$ meet at $(2,\pm i\sqrt3)$. Multiplicities: the parabola $y=x^2$ touches the line $y=0$ twice at the
origin. With these three conventions, which this chapter makes precise, B\'ezout's count is exact, and it becomes the engine of
everything else.

The second half of the chapter is about a single curve and the functions on it. In 1857 Riemann asked how many meromorphic functions a
compact Riemann surface carries with poles bounded by a given divisor $D$, and found the answer up to a correction: at least
$\deg D+1-g$, where $g$ is the number of holes of the surface. His student Gustav Roch, who died of tuberculosis at twenty-six, found the
correction term in 1865. Brill and Max Noether turned the analysis into algebra in 1874 using plane curves and B\'ezout's theorem, and
that algebra works over any algebraically closed field, in particular over $\overline{\mathbb F}_p$. This is the version proved here,
for smooth plane curves, following Fulton's book. Its payoffs come in the next three chapters: the group law on an elliptic curve
(Chapter~15), the Riemann hypothesis for curves over finite fields (Chapter~16), and, at the end of this chapter, error-correcting codes
that beat the best codes built from polynomials.

\begin{reminder}[from Volume 2]
$k[x_1,\dots,x_n]$ is a UFD (Volume~2, Theorem~6.17), and by Gauss's lemma a polynomial of $R[y]$, $R$ a UFD with fraction field $K$, that
divides another one in $K[y]$ divides it in $R[y]$ when it is primitive (Volume~2, Proposition~6.16). The Chinese remainder theorem for rings:
if $I_1,\dots,I_s$ are pairwise coprime ideals of a commutative ring $A$, then $A/(I_1\cdots I_s)\cong\prod A/I_i$ (Volume~2, Theorem~6.9, by
induction). An algebraically closed field $k$ has no proper finite extension. Throughout this chapter $k$ is an algebraically closed
field of any characteristic, except in Section~\ref{sec:goppa}.
\end{reminder}

\section{Plane curves}

An \emph{affine plane curve} is a non-constant polynomial $F\in k[x,y]$, up to a non-zero scalar factor; its points are
$V(F)=\{(a,b)\in k^2:F(a,b)=0\}$, its \emph{degree} is $\deg F$, and its \emph{components} are its irreducible factors (counted with
multiplicity). We identify a curve with its equation: different equations can have the same points ($y$ and $y^2$), and the equation is the
finer object, which carries multiplicities.

The \emph{projective plane} $\mathbb P^2$ is the set of lines through the origin of $k^3$; the line through $(a,b,c)\ne0$ is written
$(a:b:c)$. The points with $c\ne0$ form a copy of $k^2$ via $(a:b:c)\mapsto(a/c,b/c)$, and the points with $c=0$ form the \emph{line at
infinity}, one point for each direction of lines in $k^2$: the parallel lines $y=0$ and $y=1$ both pass through $(1:0:0)$. A
\emph{projective plane curve} is a non-constant homogeneous polynomial (a \emph{form}) $F\in k[X,Y,Z]$, up to a scalar; since
$F(\lambda a,\lambda b,\lambda c)=\lambda^{\deg F}F(a,b,c)$, the condition $F(P)=0$ makes sense for $P\in\mathbb P^2$. To go back and forth,
\[
F_*(x,y)=F(x,y,1),\qq f^*(X,Y,Z)=Z^{\deg f}f(X/Z,Y/Z).
\]
If $Z\nmid F$, then $(F_*)^*=F$; and $(f^*)_*=f$ always. So affine curves and projective curves without the component $Z$ correspond, and the
projective curve adds the points at infinity. Every linear change of coordinates of $k^3$ acts on $\mathbb P^2$ (a \emph{projective change of
coordinates}); it can move any line to $Z=0$ and any point to $(0:0:1)$.

Let $P=(a,b)$ be a point of the affine curve $F$. Expanding $F$ in powers of $x-a$ and $y-b$, $F=F_m+F_{m+1}+\dots$ with $F_j$ homogeneous of
degree $j$ in $x-a,y-b$ and $F_m\ne0$; the integer $m=m_P(F)$ is the \emph{multiplicity} of $P$ on $F$ (it is $0$ iff $P\notin F$). $P$ is a
\emph{smooth} (or simple) point if $m_P(F)=1$, that is if $F_x(P)$ and $F_y(P)$ are not both $0$; then the line $F_x(P)(x-a)+F_y(P)(y-b)=0$ is
the \emph{tangent line} at $P$. A point of a projective curve is smooth if it is smooth in an affine chart containing it, and a curve is
\emph{smooth} if all its points are.

\begin{lemma}\label{lem:projsmooth}
A point $P$ of a projective curve $F$ is singular iff $F_X(P)=F_Y(P)=F_Z(P)=0$. In particular this does not depend on the chart.
\end{lemma}
\begin{proof}
By symmetry let $P=(a:b:1)$ lie in the chart $Z\ne0$, where the curve is $F_*$. Then $(F_*)_x(a,b)=F_X(a,b,1)$ and $(F_*)_y(a,b)=F_Y(a,b,1)$, so
$P$ is singular iff $F_X(P)=F_Y(P)=0$. And this forces $F_Z(P)=0$, by Euler's identity $XF_X+YF_Y+ZF_Z=(\deg F)F$ (check it on monomials),
evaluated at $(a,b,1)$, where $F$ vanishes. This holds in every characteristic.
\end{proof}

\begin{example}[three cubics]\label{exa:threecubics}
The cubic $y^2=x^2(x+1)$ has a \emph{node} at the origin: its lowest terms are $y^2-x^2=(y-x)(y+x)$, two tangent lines. The cubic $y^2=x^3$ has a
\emph{cusp}: lowest term $y^2$, one tangent counted twice. The cubic $y^2=x^3-x$ is smooth: $F=y^2-x^3+x$ has $F_y=2y$, $F_x=1-3x^2$, which do not
vanish together on the curve if $\operatorname{char}k\ne2$ (if $y=0$ then $x\in\{0,\pm1\}$ and $1-3x^2\in\{1,-2\}$); at infinity
$Y^2Z-X^3+XZ^2$ has the single point $(0:1:0)$, where $F_Z=Y^2+2XZ=1\ne0$. Figure~\ref{fig:threecubics} shows the real points. The first two are
\emph{rational}: the line $y=tx$ through the singular point meets the curve in one more point, $(t^2-1,t(t^2-1))$, respectively $(t^2,t^3)$, and
this parametrises the curve by $t$. The third one admits no such parametrisation (Corollary~\ref{cor:genus1} will say in what sense).
\end{example}

\begin{figure}[ht]
\centering\fig{ch11-cubics}
\caption{The real points of a nodal cubic $y^2=x^2(x+1)$, a cuspidal cubic $y^2=x^3$, and a smooth cubic $y^2=x^3-x$, which has two real
components; the line through the singular point used to parametrise the first two is dashed.}
\label{fig:threecubics}
\end{figure}

The first thing to know about two curves is that they meet in finitely many points, with a finite amount of algebra at each.

\begin{lemma}\label{lem:finite}
Let $F,G\in k[x,y]$ have no common factor. Then $(F,G)$ contains a non-zero polynomial in $x$ alone and a non-zero polynomial in $y$ alone.
Hence $V(F)\cap V(G)$ is finite and $k[x,y]/(F,G)$ is a finite-dimensional $k$-vector space.
\end{lemma}
\begin{proof}
View $F,G$ in $R[y]$ with $R=k[x]$, and in $K[y]$ with $K=k(x)$. They have no common factor of positive degree in $K[y]$: such a factor can
be scaled to a primitive polynomial of $R[y]$, which would divide both in $R[y]$ (Volume~2, Proposition~6.16). If one of $F,G$ is a non-zero
element of $R$ there is nothing to prove; otherwise, as $K[y]$ is a PID, $AF+BG=1$ for some $A,B\in K[y]$. Multiplying by a common denominator
$D(x)\in R$ gives $aF+bG=D(x)\ne0$ with $a,b\in k[x,y]$. By symmetry $(F,G)$ also contains some $E(y)\ne0$. The points of $V(F)\cap V(G)$ have their
$x$-coordinate among the roots of $D$ and their $y$-coordinate among those of $E$, so they are finitely many; and modulo $(D(x),E(y))$ every
polynomial reduces to a combination of the $x^iy^j$ with $i<\deg D$, $j<\deg E$.
\end{proof}

For a point $P$ of $k^2$, let $\mathcal O_P=\mathcal O_P(\mathbb A^2)$ be the ring of rational functions $f/g$ with $f,g\in k[x,y]$ and $g(P)\ne0$:
the functions defined near $P$. It is a local ring: its non-units are the functions vanishing at $P$, and they form its unique maximal ideal
$\mathfrak m_P$. Ideals of $k[x,y]$ generate ideals of $\mathcal O_P$, written $I\mathcal O_P$.

\begin{lemma}[local decomposition]\label{lem:localdecomp}
Let $I$ be an ideal of $k[x,y]$ with $A=k[x,y]/I$ finite-dimensional. Then $V(I)=\{P_1,\dots,P_s\}$ is finite and
\[
k[x,y]/I\;\cong\;\prod_{i=1}^s\mathcal O_{P_i}/I\mathcal O_{P_i},
\]
the map being induced by the inclusions $k[x,y]\subseteq\mathcal O_{P_i}$. In particular $h\in I$ iff $h\in I\mathcal O_P$ for every $P\in V(I)$.
\end{lemma}
\begin{proof}
\emph{Maximal ideals.} Let $M$ be a maximal ideal of $A$. The field $A/M$ is finite-dimensional over $k$, so it is $k$; if $x\mapsto a$ and $y\mapsto b$
in it, then $M$ contains the images of $x-a$ and $y-b$, and since $(x-a,y-b)$ is already maximal, $M$ is its image. So the maximal ideals of $A$
are the images $M_i$ of $\mathfrak m_{P}\cap k[x,y]=(x-a,y-b)$ for the points $P=(a,b)$ with $I\subseteq(x-a,y-b)$, that is $P\in V(I)$. Distinct
maximal ideals are pairwise coprime, so by the Chinese remainder theorem $A\to\prod A/M_i$ is onto, and there are at most $\dim A$ of them.

\emph{Nilpotence.} Every element $u$ of $A$ outside all the $M_i$ is a unit: otherwise $uA$ is a proper ideal, contained in a maximal one. Let
$a\in J=\bigcap M_i$, and let $\mu(T)=T^j\nu(T)$, $\nu(0)\ne0$, be the minimal polynomial of the linear map ``multiplication by $a$''. Then
$\nu(a)\equiv\nu(0)\not\equiv0$ modulo each $M_i$, so $\nu(a)$ is a unit, and $a^j\nu(a)=0$ gives $a^j=0$. So $J$ is spanned by finitely many
nilpotent elements, say $N$ of them with $r$-th powers $0$, and any product of $Nr$ elements of $J$ vanishes: $J^{Nr}=0$. Put $R=Nr$.

\emph{Splitting.} The ideals $M_i^R$ are pairwise coprime (raise $1=u+v$, $u\in M_i$, $v\in M_j$, to the power $2R$), and their product is contained
in $J^R=0$. By the Chinese remainder theorem, $A\cong\prod_iA/M_i^R$. Each factor $A_i=A/M_i^R$ is local: its only maximal ideal is the image of
$M_i$, and every element outside it is a unit. Finally $A_i\cong\mathcal O_{P_i}/I\mathcal O_{P_i}$: the map $k[x,y]\to A_i$ sends every
$g$ with $g(P_i)\ne0$ to a unit, so it extends to $\mathcal O_{P_i}\to A_i$, which kills $I\mathcal O_{P_i}$; it is onto. Conversely, let
$e_i\in A$ be the idempotent that is $1$ in $A_i$ and $0$ in the other factors, and lift it to $\tilde e\in k[x,y]$; then $\tilde e(P_i)=1$, because $e_i-1$ maps to $0$ in $A_i$ and so lies in the image of $M_i$. If $h\in k[x,y]$ maps to $0$ in $A_i$, then $he_i=0$ in $A$, so $h\tilde e\in I$, and
$h=h\tilde e/\tilde e\in I\mathcal O_{P_i}$ since $\tilde e$ is a unit of $\mathcal O_{P_i}$. Every element of $\mathcal O_{P_i}$ is $h/g$, so
the kernel of $\mathcal O_{P_i}\to A_i$ is exactly $I\mathcal O_{P_i}$.
\end{proof}

Lemmas~\ref{lem:finite} and~\ref{lem:localdecomp} together are the two-variable case of Hilbert's Nullstellensatz, which is all this chapter
needs.

\section{Intersection numbers}

\begin{definition}
Let $F,G\in k[x,y]$ have no common factor and let $P\in k^2$. The \emph{intersection number} of $F$ and $G$ at $P$ is
\[
I(P,F\cap G)=\dim_k\mathcal O_P/(F,G)\mathcal O_P.
\]
\end{definition}

It is finite, being the dimension of a factor in Lemma~\ref{lem:localdecomp} (with $I=(F,G)$, which is allowed by Lemma~\ref{lem:finite}), and
the same lemma says
\begin{equation}\label{eq:sumI}
\sum_{P}I(P,F\cap G)=\dim_kk[x,y]/(F,G).
\end{equation}
Two examples. For the axes, $\mathcal O_P/(x,y)=k$ at $P=(0,0)$, so $I=1$. For the parabola and its tangent, $(y-x^2,y)=(y,x^2)$, and
$\mathcal O_P/(y,x^2)$ has basis $1,x$: $I=2$, the tangency counted twice.

\begin{proposition}\label{prop:intprops}
Let $F,G,H\in k[x,y]$, with $F$ prime to $G$ and to $H$, and $P\in k^2$.
\begin{enumerate}[label=(\alph*)]
\item $I(P,F\cap G)=I(P,G\cap F)$; it is $0$ iff $P\notin V(F)\cap V(G)$; it is unchanged by an affine change of coordinates.
\item $I(P,F\cap G)=I(P,F\cap(G+AF))$ for every $A\in k[x,y]$.
\item $I(P,F\cap GH)=I(P,F\cap G)+I(P,F\cap H)$.
\end{enumerate}
\end{proposition}
\begin{proof}
(a) The ideal $(F,G)$ is symmetric. If $F(P)\ne0$ then $F$ is a unit of $\mathcal O_P$ and the quotient is $0$; if $F(P)=G(P)=0$ then
$(F,G)\mathcal O_P\subseteq\mathfrak m_P$ and the quotient maps onto $\mathcal O_P/\mathfrak m_P=k$. An affine change of coordinates induces
an isomorphism of the local rings carrying one ideal to the other. (b) The ideals $(F,G)$ and $(F,G+AF)$ are equal.

(c) Write $\mathcal O=\mathcal O_P$. Multiplication by $G$ gives a map $\mathcal O/(F,H)\to\mathcal O/(F,GH)$, and $\mathcal O/(F,GH)\to\mathcal O/(F,G)$ is
the natural projection. The sequence
\[
0\to\mathcal O/(F,H)\xrightarrow{\;\cdot G\;}\mathcal O/(F,GH)\to\mathcal O/(F,G)\to0
\]
is exact: the projection is onto, the image of $\cdot G$ is $(F,G)/(F,GH)$, which is the kernel of the projection, and $\cdot G$ is injective.
For the last point, suppose $Gz\in(F,GH)\mathcal O$; multiplying by a polynomial that does not vanish at $P$, we may take $z,u,v\in k[x,y]$ with
$sGz=uF+vGH$, $s(P)\ne0$. Then $G(sz-vH)=uF$ in the UFD $k[x,y]$, and $F$ is prime to $G$, so $F\mid sz-vH$ and $sz\in(F,H)$, hence
$z\in(F,H)\mathcal O$. Dimensions add along an exact sequence.
\end{proof}

\subsection*{Smooth points}

When $P$ is a smooth point of $F$, intersection numbers become orders of vanishing. A smooth point lies on only one component of $F$
(a point on two components, or on a repeated one, has multiplicity $\ge2$), and the other components do not vanish at $P$, so they are units of
$\mathcal O_P$. Hence the \emph{local ring of $F$ at $P$},
\[
\mathcal O_P(F)=\mathcal O_P/F\mathcal O_P,
\]
equals $\mathcal O_P/F_1\mathcal O_P$ for the component $F_1$ through $P$. It is a domain: $F_1$ is irreducible, hence prime in the UFD
$k[x,y]$, and if $F_1\mid ab$ in $\mathcal O_P$, clearing denominators, $F_1$ divides one of the numerators.

\begin{proposition}\label{prop:smoothdvr}
Let $P$ be a smooth point of $F$, and let $t$ be (the image of) the equation of any line through $P$ other than the tangent. Then every
non-zero $z\in\mathcal O_P(F)$ can be written uniquely $z=ut^r$ with $u$ a unit and $r\ge0$. The integer $r=\operatorname{ord}_P(z)$ satisfies
$\operatorname{ord}_P(zz')=\operatorname{ord}_P(z)+\operatorname{ord}_P(z')$, and for every $G$ prime to $F$,
\[
I(P,F\cap G)=\operatorname{ord}_P(G).
\]
\end{proposition}
\begin{proof}
Change coordinates so that $P=(0,0)$, the line is $t=x$, and the tangent is $y=0$; scaling, $F=y+(\text{terms of degree}\ge2)$. Every monomial of
degree $\ge2$ is divisible by $y$ or is a power $x^j$, $j\ge2$, so $F=y(1+G_1)+x^2H(x)$ with $G_1(P)=0$. In $\mathcal O_P(F)$, $y=-x^2H/(1+G_1)$ is a
multiple of $x$. So the maximal ideal of $\mathcal O_P(F)$, generated by $x$ and $y$, is $(t)$.

Let $z\ne0$, written $z=g/h$ with $g$ not divisible by $F_1$. If $z\in(t^n)$ then $(g)\subseteq(t^n)$ in $\mathcal O_P(F)$, so
$\dim\mathcal O_P(F)/(t^n)\le\dim\mathcal O_P(F)/(g)=I(P,F_1\cap g)<\infty$. But the chain $\mathcal O_P(F)\supsetneq(t)\supsetneq\dots\supsetneq(t^n)$ is
strict ($t^i=t^{i+1}w$ would give $t^i(1-tw)=0$ with $1-tw$ a unit, so $t^i=0$, absurd in a domain), so $\dim\mathcal O_P(F)/(t^n)\ge n$. Hence
there is a largest $r$ with $z\in(t^r)$; then $z=ut^r$ with $u\notin(t)$, a unit. If $ut^r=u't^{r'}$ with $r\le r'$, then $u=u't^{r'-r}$ and
$r=r'$ since $u$ is a unit. Additivity follows from uniqueness. Finally $\mathcal O_P/(F,G)=\mathcal O_P(F)/(G)=\mathcal O_P(F)/(t^r)$ with
$r=\operatorname{ord}_P(G)$, and $(t^i)/(t^{i+1})$ is spanned by $t^i$ and non-zero, for $i<r$: the dimension is $r$.
\end{proof}

So $\mathcal O_P(F)$ is a \emph{discrete valuation ring}, the curve analogue of $\Z_{(p)}$ or of $\Z_p$ in Chapter~8, with uniformiser $t$.
On the field of fractions $\operatorname{ord}_P(g/h)=\operatorname{ord}_P(g)-\operatorname{ord}_P(h)$.

\subsection*{Projective intersection numbers}

For a point $P$ of $\mathbb P^2$, let $\mathcal O_P(\mathbb P^2)$ be the ring of quotients $G/H$ of forms of the same degree with $H(P)\ne0$; these are
honest functions near $P$. If $P$ lies in the chart $Z\ne0$, dehomogenisation $G/H\mapsto G_*/H_*$ identifies $\mathcal O_P(\mathbb P^2)$ with
$\mathcal O_P(\mathbb A^2)$, with inverse $f/g\mapsto Z^{\deg g-\deg f}f^*/g^*$ (multiply the numerator or the denominator by a power of $Z$ to equalise
degrees). For forms $F,G$ without common component, and any linear form $L$ with $L(P)\ne0$, put
\[
I(P,F\cap G)=\dim_k\mathcal O_P(\mathbb P^2)\big/\bigl(F/L^{\deg F},\,G/L^{\deg G}\bigr).
\]
Changing $L$ multiplies the generators by the units $(L/L')^{\deg F}$ and $(L/L')^{\deg G}$, so this is independent of $L$; with $L=Z$ it is the
affine intersection number of $F_*$ and $G_*$. Proposition~\ref{prop:intprops} and Proposition~\ref{prop:smoothdvr} carry over, with $\mathcal O_P(F)=\mathcal O_P(\mathbb P^2)/(F/L^{\deg F})$, and they are now invariant under projective changes of coordinates. For a
form $G$ we write $\operatorname{ord}_P(G)$ for $\operatorname{ord}_P(G/L^{\deg G})$.

\section{B\'ezout's theorem}

\begin{theorem}[B\'ezout]\label{thm:bezout}
Let $F$ and $G$ be projective plane curves of degrees $m$ and $n$ with no common component. Then
\[
\sum_{P\in\mathbb P^2}I(P,F\cap G)=mn.
\]
\end{theorem}
\begin{proof}
\emph{Coordinates.} $F$ and $G$ meet in finitely many points (Lemma~\ref{lem:finite} in each of the three charts). As $k$ is infinite, some line
avoids them all; make it $Z=0$. Then all the points lie in the affine chart, and by \eqref{eq:sumI} the sum is $\dim_k\Gamma_*$, with
$\Gamma_*=k[x,y]/(F_*,G_*)$. We show $\dim\Gamma_*=mn$.

\emph{Graded pieces.} Let $R=k[X,Y,Z]$, $R_d$ its forms of degree $d$ (with $0$), $\dim R_d=\binom{d+2}2$, and $\Gamma=R/(F,G)$ with graded pieces
$\Gamma_d=R_d/(FR_{d-m}+GR_{d-n})$. The sequence
\[
0\to R_{d-m-n}\xrightarrow{\;C\mapsto(GC,-FC)\;}R_{d-m}\times R_{d-n}\xrightarrow{\;(A,B)\mapsto AF+BG\;}R_d\to\Gamma_d\to0
\]
is exact. The only point to check is in the middle: if $AF+BG=0$, then $F\mid BG$, so $F\mid B$ ($R$ is a UFD and $F,G$ have no common factor),
$B=-FC$, and then $A=GC$. Taking alternating sums of dimensions, for $d\ge m+n$,
\[
\dim\Gamma_d=\tbinom{d+2}2-\tbinom{d-m+2}2-\tbinom{d-n+2}2+\tbinom{d-m-n+2}2=mn,
\]
a polynomial identity in $d$ that one checks by expanding (the $d^2$ and $d$ terms cancel).

\emph{$Z$ is a non-zero-divisor on $\Gamma$.} Suppose $ZH=AF+BG$. Set $Z=0$: writing $A_0(X,Y)=A(X,Y,0)$ and so on, $A_0F_0+B_0G_0=0$. The binary
forms $F_0,G_0$ are non-zero (if, say, $F_0=0$, every zero of $G$ on $Z=0$ would be a common point there) and have no common factor: a binary form over $k$ factors into linear forms $bX-aY$, and a common factor would give a common zero
$(a:b:0)$ of $F$ and $G$ on the line $Z=0$. So, as before, $B_0=-F_0C$, $A_0=G_0C$ for a form $C\in k[X,Y]$. Then $A-GC$ and $B+FC$ vanish on $Z=0$,
that is $A-GC=ZA'$, $B+FC=ZB'$, and $ZH=(A-GC)F+(B+FC)G=Z(A'F+B'G)$. Cancelling $Z$, $H\in(F,G)$.

\emph{Conclusion.} For $d\ge m+n$, multiplication by $Z$ is an injective map $\Gamma_d\to\Gamma_{d+1}$ between spaces of dimension $mn$, hence
bijective. Consider $\alpha:\Gamma_d\to\Gamma_*$, $H\mapsto H_*$. It is injective: if $H_*=aF_*+bG_*$, homogenising gives $Z^rH=AF+BG$ for some $r$,
and $H\in(F,G)$ because $Z$ is a non-zero-divisor. It is onto: given $h\in k[x,y]$, take $N\ge\deg h$; then $Z^{d+N-\deg h}h^*\in R_{d+N}$, and since $Z^N:\Gamma_d\to\Gamma_{d+N}$ is
bijective, $Z^{d+N-\deg h}h^*\equiv Z^NH\pmod{(F,G)}$ for some $H\in R_d$; dehomogenising, $h\equiv H_*$ modulo $(F_*,G_*)$. So $\dim\Gamma_*=\dim\Gamma_d=mn$.
\end{proof}

\begin{think}
The lines $y=0$ and $y=1$ are parallel, yet B\'ezout says they meet once. Where? And the circles $x^2+y^2=1$ and $(x-1)^2+y^2=1$ meet in
two real points: where are the other two? (The second answer is Exercise~\ref{ex:circles}.)
\end{think}

\begin{corollary}\label{cor:bezoutcons}
\begin{enumerate}[label=(\alph*)]
\item Any two projective plane curves meet.
\item A line not contained in a curve of degree $n$ meets it in exactly $n$ points counted with multiplicity.
\item A smooth projective plane curve is irreducible.
\end{enumerate}
\end{corollary}
\begin{proof}
(a) and (b) are Theorem~\ref{thm:bezout} (with common components, every point of the common one is in both). (c) If $F=GH$ with $G,H$ non-constant, take
$P\in G\cap H$, by (a). Then $F_X=G_XH+GH_X$ vanishes at $P$, and so do $F_Y,F_Z$: $P$ is singular by Lemma~\ref{lem:projsmooth}.
\end{proof}

Figure~\ref{fig:bezout} shows the count for two parabolas, $y=x^2$ and $y=2x^2-1$. In the affine plane they meet at $(\pm1,1)$, transversally:
two points. B\'ezout promises four. Homogenised, $YZ=X^2$ and $YZ=2X^2-Z^2$ both pass through $(0:1:0)$. In the chart $Y=1$, with coordinates
$x=X/Y$, $z=Z/Y$, the curves are $z=x^2$ and $z=2x^2-z^2$; substituting, $(z-x^2,\,x^2-z^2)=(z-x^2,\,x^2-x^4)=(z-x^2,\,x^2(1-x^2))$, and since $1-x^2$ is a unit
at the origin, the local ring modulo this ideal is $\mathcal O/(z-x^2,x^2)$, of dimension $2$. The parabolas are tangent to each other, and to the
line at infinity, at the point at infinity of the $y$-axis: $2+1+1=4$.

\begin{figure}[ht]
\centering\fig{ch11-bezout}
\caption{The parabolas $y=x^2$ and $y=2x^2-1$ meet in two points of the affine plane; the other two of B\'ezout's four are at the point at infinity
$(0:1:0)$ of the vertical direction, where the curves are tangent.}
\label{fig:bezout}
\end{figure}

\section{Max Noether's theorem}

When does a third curve $H$ belong to the ideal $(F,G)$, that is $H=AF+BG$ with forms $A,B$? Necessarily locally at each point of $F\cap G$.
Max Noether's theorem says that this is enough.

\begin{theorem}[Noether's $AF+BG$ theorem]\label{thm:noether}
Let $F,G,H$ be forms, $F$ and $G$ without common component. Suppose that at every $P\in F\cap G$, the germ of $H$ lies in the ideal generated by
those of $F$ and $G$ in $\mathcal O_P(\mathbb P^2)$ (Noether's conditions). Then $H=AF+BG$ with forms $A,B$ of degrees $\deg H-\deg F$ and
$\deg H-\deg G$ (or zero).
\end{theorem}
\begin{proof}
As in the proof of Theorem~\ref{thm:bezout}, choose $Z=0$ avoiding $F\cap G$. In the chart, Noether's conditions say $H_*\in(F_*,G_*)\mathcal O_P$ for
every $P\in V(F_*,G_*)$, so $H_*\in(F_*,G_*)$ by Lemma~\ref{lem:localdecomp}: $H_*=aF_*+bG_*$. Homogenising, $Z^rH=A'F+B'G$, and since $Z$ is a non-zero-divisor
modulo $(F,G)$ (shown in the proof of Theorem~\ref{thm:bezout}), $H=A''F+B''G$. Keep only the terms of degree $\deg H$: $H=AF+BG$ with $A,B$ the homogeneous
parts of $A'',B''$ of the right degrees.
\end{proof}

At smooth points the conditions are a simple inequality.

\begin{corollary}\label{cor:noethersmooth}
Let $F$ be a curve, $G$ and $H$ forms prime to $F$, and suppose every point of $F\cap G$ is a smooth point of $F$. If
$I(P,F\cap H)\ge I(P,F\cap G)$ for every $P$, then $H=AF+BG$ with forms $A,B$.
\end{corollary}
\begin{proof}
At $P\in F\cap G$, the ring $\mathcal O_P(F)$ is a discrete valuation ring (Proposition~\ref{prop:smoothdvr}), where $G=ut^r$ with
$r=I(P,F\cap G)$ and $H=u't^{r'}$ with $r'\ge r$. So $H=(u'u^{-1}t^{r'-r})G$ in $\mathcal O_P(F)$, that is $H\in(F,G)$ in
$\mathcal O_P(\mathbb P^2)$. Apply Theorem~\ref{thm:noether}.
\end{proof}

\section{Divisors on a smooth curve}

From now on $C$ is a smooth projective plane curve, given by a form $F$ of degree $n$; it is irreducible by Corollary~\ref{cor:bezoutcons}.
Its \emph{function field} $k(C)$ is the set of quotients $G/H$ of forms of the same degree with $F\nmid H$, two quotients $G/H$ and $G'/H'$ being
equal when $F\mid GH'-G'H$. It is a field (because $F$ is prime), and it is the field of fractions of every $\mathcal O_P(C)$. Each point
$P$ of $C$ gives a valuation $\operatorname{ord}_P$ on $k(C)^\times$ (Proposition~\ref{prop:smoothdvr}).

A \emph{divisor} is a formal finite sum $D=\sum_Pn_PP$ of points of $C$ with integer coefficients; its \emph{degree} is $\sum n_P$, and $D\ge D'$ means
$n_P\ge n'_P$ for all $P$. A divisor $D\ge0$ is \emph{effective}. For a form $G$ prime to $F$ and for $z\in k(C)^\times$ set
\[
\operatorname{div}(G)=\sum_PI(P,F\cap G)\,P,\qquad\operatorname{div}(z)=\sum_P\operatorname{ord}_P(z)\,P.
\]
By Proposition~\ref{prop:smoothdvr}, $\operatorname{div}(G/H)=\operatorname{div}(G)-\operatorname{div}(H)$; it is a finite sum, $\operatorname{div}(GG')=\operatorname{div}(G)+\operatorname{div}(G')$, and by
B\'ezout $\deg\operatorname{div}(G)=n\deg G$. Divisors $\operatorname{div}(z)$ are \emph{principal}, and $D\equiv D'$ (\emph{linear equivalence}) means that
$D-D'$ is principal. We fix a line $L$ with $L\nmid F$ and write
\[
E=\operatorname{div}(L),\qquad E_d=dE=\operatorname{div}(L^d),
\]
an effective divisor of degree $n$; changing $L$ changes $E$ only up to linear equivalence ($E-E'=\operatorname{div}(L/L')$).

\begin{proposition}\label{prop:divbasic}
\begin{enumerate}[label=(\alph*)]
\item Principal divisors have degree $0$.
\item If $z\in k(C)^\times$ has no pole ($\operatorname{div}(z)\ge0$), then $z$ is a constant.
\end{enumerate}
\end{proposition}
\begin{proof}
(a) $\deg\operatorname{div}(G/H)=n\deg G-n\deg H=0$. (b) Write $z=G/H$. Then $\operatorname{div}(G)\ge\operatorname{div}(H)$, so by Corollary~\ref{cor:noethersmooth}
$G=AF+BH$ with $B$ a form of degree $\deg G-\deg H=0$, a constant $b$; in $k(C)$, $z=G/H=b$.
\end{proof}

For a divisor $D$, the \emph{Riemann--Roch space} is
\[
L(D)=\{z\in k(C)^\times:\operatorname{div}(z)+D\ge0\}\cup\{0\},
\]
the functions whose poles are bounded by the positive part of $D$ and which vanish at least as prescribed by its negative part. It is a vector
space, because $\operatorname{ord}_P(z+z')\ge\min(\operatorname{ord}_Pz,\operatorname{ord}_Pz')$. Let $\ell(D)=\dim_kL(D)$.

\begin{proposition}\label{prop:ellbasic}
\begin{enumerate}[label=(\alph*)]
\item If $D\equiv D'$, then $\ell(D)=\ell(D')$. $\ell(0)=1$, and $\ell(D)=0$ when $\deg D<0$.
\item $\ell(D+P)\le\ell(D)+1$ for every point $P$. Hence $\ell(D)\le\deg D+1$ for $D\ge0$, and $\ell(D)$ is always finite.
\end{enumerate}
\end{proposition}
\begin{proof}
(a) If $D'=D+\operatorname{div}(w)$, then $z\mapsto z/w$ maps $L(D)$ onto $L(D')$. $L(0)$ is the constants by Proposition~\ref{prop:divbasic}(b). If
$z\in L(D)$ is non-zero then $\deg D=\deg(\operatorname{div}(z)+D)\ge0$.

(b) Let $m$ be the coefficient of $P$ in $D$, and $t$ a uniformiser at $P$. For $z\in L(D+P)$, $t^{m+1}z\in\mathcal O_P(C)$, and $z\mapsto(t^{m+1}z)(P)$ is a linear form
on $L(D+P)$ whose kernel is $L(D)$. Induction from $\ell(0)=1$ gives $\ell(D)\le\deg D+1$ for $D\ge0$, and in general $L(D)\subseteq L(D^+)$, where $D^+\ge D$ keeps
only the positive coefficients.
\end{proof}

The next lemma says that the functions in $L(E_d-D)$ are exactly the curves of degree $d$ ``through $D$''.

\begin{lemma}\label{lem:noetherrep}
Let $D\ge0$ and $d\ge0$. Then $L(E_d-D)=\{G/L^d:\ G\in R_d,\ G=0\text{ or }\operatorname{div}(G)\ge D\}$. Two forms give the same function iff they
are congruent modulo $F$.
\end{lemma}
\begin{proof}
If $\operatorname{div}(G)\ge D$ then $\operatorname{div}(G/L^d)+E_d-D=\operatorname{div}(G)-D\ge0$. Conversely let $z\in L(E_d-D)$, $z=G_0/H_0$. Then
$\operatorname{div}(G_0L^d)=\operatorname{div}(z)+E_d+\operatorname{div}(H_0)\ge D+\operatorname{div}(H_0)\ge\operatorname{div}(H_0)$, so by
Corollary~\ref{cor:noethersmooth}, $G_0L^d=AF+GH_0$ with $G$ a form of degree $d$. In $k(C)$, $z=G_0/H_0=G/L^d$, and
$\operatorname{div}(G)=\operatorname{div}(z)+E_d\ge D$. The last sentence is the definition of equality in $k(C)$.
\end{proof}

The following theorem, Noether's \emph{Restsatz}, lets one slide part of an intersection along a linear equivalence while keeping the rest fixed.

\begin{theorem}[residue theorem]\label{thm:residue}
Let $D,D'\ge0$ with $D\equiv D'$, and let $G$ be a form with $\operatorname{div}(G)=D+A$, $A\ge0$. Then there is a form $G'$ of the same degree with
$\operatorname{div}(G')=D'+A$.
\end{theorem}
\begin{proof}
Write $D'-D=\operatorname{div}(H/H')$ with forms $H,H'$ of the same degree. Then
$\operatorname{div}(GH)=D+A+\operatorname{div}(H)=D'+A+\operatorname{div}(H')\ge\operatorname{div}(H')$, so by Corollary~\ref{cor:noethersmooth},
$GH=BF+G'H'$ with $G'$ a form of degree $\deg G$. Then
$\operatorname{div}(G')=\operatorname{div}(GH)-\operatorname{div}(H')=D'+A$.
\end{proof}

\section{The Riemann--Roch theorem}

\begin{definition}
The \emph{genus} of the smooth plane curve $C$ of degree $n$ is $g=\frac{(n-1)(n-2)}2$. Its \emph{canonical divisor} is $W=E_{n-3}=(n-3)E$, of
degree $n(n-3)=2g-2$.
\end{definition}

Lines and conics have genus $0$, cubics genus $1$, quartics genus $3$. Over $\C$, a smooth plane curve is a compact Riemann surface, and $g$ is its
number of holes (a cubic is a torus, Chapter~15). The canonical divisor is the divisor of the differential form $\omega=\mathrm dx/F_y$ (in the
chart $Z\ne0$, with $F$ standing for $F_*$), whose zeros and poles all lie at infinity; we shall not need differentials, and take the definition as
it stands. What matters is its class, and $\ell(W)=g$ by Lemma~\ref{lem:noetherrep}: $L(E_{n-3})$ is the space of all forms of degree $n-3$, of
dimension $\binom{n-1}2=g$ (none is divisible by $F$).

The proof of the Riemann--Roch theorem has three steps: count forms to get Riemann's inequality, then show, through B\'ezout's theorem on a line, that the
two sides of the formula never grow at the same time, then put the two together.

\begin{proposition}[counting forms]\label{prop:countforms}
Let $V_d=R_d/FR_{d-n}$, the forms of degree $d$ modulo $F$. For $d\ge n-2$, $\dim V_d=dn+1-g$, and $\ell(E_d)=dn+1-g$.
\end{proposition}
\begin{proof}
Multiplication by $F$ is injective, so $\dim V_d=\binom{d+2}2-\binom{d-n+2}2$, and $\binom{d-n+2}2=\frac{(d-n+2)(d-n+1)}2$ holds for $d-n\ge-2$ (both sides vanish at $-1$ and $-2$).
Expanding, $\frac{(d+2)(d+1)-(d-n+2)(d-n+1)}2=dn+1-\frac{(n-1)(n-2)}2$. And $L(E_d)\cong V_d$ by Lemma~\ref{lem:noetherrep} with $D=0$.
\end{proof}

\begin{proposition}[Riemann's inequality]\label{prop:riemann}
For every divisor $D$, $\ell(D)\ge\deg D+1-g$.
\end{proposition}
\begin{proof}
\emph{Effective divisors with a residual.} Let $D_0\ge0$ of degree $r$ and $d\ge n-2$. For each $P$ in the support of $D_0$, with coefficient $m_P$, the map
$G\mapsto G/L_P^d\bmod\mathfrak m_P^{m_P}$ (with $L_P$ a line not through $P$) sends $V_d$ to $\mathcal O_P(C)/\mathfrak m_P^{m_P}$, of dimension $m_P$; the
kernel of $V_d\to\prod_P\mathcal O_P(C)/\mathfrak m_P^{m_P}$ is the set of $G$ with $\operatorname{div}(G)\ge D_0$, of dimension at least
$\dim V_d-r$. By Lemma~\ref{lem:noetherrep}, $\ell(E_d-D_0)\ge dn+1-g-r=\deg(E_d-D_0)+1-g$.

\emph{Effective divisors.} Let $D\ge0$. Choose $d\ge n-2$ with $dn+1-g>\deg D$; by the count just made there is a form $G$ of degree $d$, not divisible by $F$,
with $\operatorname{div}(G)\ge D$. Write $\operatorname{div}(G)=D+D_0$, $D_0\ge0$. Since $\operatorname{div}(G)-E_d=\operatorname{div}(G/L^d)$, $D\equiv E_d-D_0$, and
$\ell(D)=\ell(E_d-D_0)\ge\deg D+1-g$.

\emph{All divisors.} Write $D=D^+-D^-$ with $D^\pm\ge0$. Applying Proposition~\ref{prop:ellbasic}(b) $\deg D^-$ times,
$\ell(D)\ge\ell(D^+)-\deg D^-\ge\deg D^++1-g-\deg D^-$.
\end{proof}

Now for the correction term. The key is to compare a curve's behaviour along $C$ and along a line.

\begin{lemma}[truncated intersection numbers]\label{lem:truncate}
Let $Q$ be a point of $C$, $M$ a line through $Q$ not contained in $C$, and $e=I(Q,C\cap M)$. For every form $H$ prime to $F$ and to $M$,
\[
\min\bigl(I(Q,C\cap H),e\bigr)=\min\bigl(I(Q,M\cap H),e\bigr).
\]
\end{lemma}
\begin{proof}
Let $\mathcal O=\mathcal O_Q(\mathbb P^2)$ and $B=\mathcal O/(F,M)$, of dimension $e$ (dehomogenised generators understood). $B$ is a quotient of
$\mathcal O_Q(C)$, by the ideal generated by $M$, which is $(t^e)$ for a uniformiser $t$ of $C$ at $Q$ (Proposition~\ref{prop:smoothdvr}); so
$B=\mathcal O_Q(C)/(t^e)$. Symmetrically, $Q$ is a smooth point of the line $M$, and $B=\mathcal O_Q(M)/(s^e)$ for a uniformiser $s$ of $M$. In both
descriptions the maximal ideal $\mathfrak n$ of $B$ is the image of $\mathfrak m_Q$, and for $i\le e$, $\mathfrak n^i$ is the image of $(t^i)$ and of
$(s^i)$. So, for $h\in\mathcal O$, the largest $i\le e$ with $\bar h\in\mathfrak n^i$ is $\min(\operatorname{ord}^C_Q(h),e)$ by the first description and
$\min(\operatorname{ord}^M_Q(h),e)$ by the second. Proposition~\ref{prop:smoothdvr} turns these orders into the two intersection numbers.
\end{proof}

\begin{lemma}[Noether's reduction lemma]\label{lem:reduction}
Let $D\ge0$ and $P\in C$. If $\ell(W-D-P)\ne\ell(W-D)$, then $\ell(D+P)=\ell(D)$.
\end{lemma}
\begin{proof}
We may assume $n\ge3$: otherwise $\deg(W-D)<0$ and $\ell(W-D)=0=\ell(W-D-P)$. Take $z_0\in L(W-D)\setminus L(W-D-P)$. By Lemma~\ref{lem:noetherrep},
$z_0=H/L^{n-3}$ with $H$ of degree $n-3$ and $\operatorname{div}(H)=\operatorname{div}(z_0)+W=D+A$, $A\ge0$; and $z_0\notin L(W-D-P)$ says that $P$ does
not occur in $A$.

Suppose, for a contradiction, that some $z\in L(D+P)$ is not in $L(D)$. Then $D'=D+P+\operatorname{div}(z)$ is effective, $D'\equiv D+P$, and $P$ does not
occur in $D'$ (the coefficient of $P$ is exactly $0$, since $z\notin L(D)$). Let $M$ be any line through $P$ not contained in $C$, and
$\operatorname{div}(M)=P+B$, $B\ge0$. The form $HM$ of degree $n-2$ has $\operatorname{div}(HM)=(D+P)+(A+B)$, so by the residue theorem there is a form
$H'$ of degree $n-2$ with
\[
\operatorname{div}(H')=D'+A+B.
\]
Compare $H'$ with $M$ at each $Q\in C\cap M$, with $e_Q=I(Q,C\cap M)$. If $Q\ne P$, the coefficient of $Q$ in $B$ is $e_Q$, so
$I(Q,C\cap H')\ge e_Q$. At $P$, the coefficient of $P$ in $B$ is $e_P-1$, and $P$ occurs neither in $D'$ nor in $A$, so $I(P,C\cap H')=e_P-1$. If $M\mid H'$,
then $I(P,C\cap H')\ge I(P,C\cap M)=e_P$, a contradiction. Otherwise Lemma~\ref{lem:truncate} gives $I(Q,M\cap H')\ge e_Q$ for $Q\ne P$ and
$I(P,M\cap H')=e_P-1$, so, by B\'ezout for $M$ and $H'$, $n-2=\sum_QI(Q,M\cap H')\ge\sum_Qe_Q-1=n-1$: a contradiction again.
\end{proof}

\begin{theorem}[Riemann--Roch]\label{thm:rr}
For every divisor $D$ on a smooth projective plane curve $C$ of genus $g$ and canonical divisor $W$,
\[
\ell(D)-\ell(W-D)=\deg D+1-g.
\]
\end{theorem}
\begin{proof}
Let $\Delta(D)=\ell(D)-\ell(W-D)-\deg D-1+g$; it only depends on the class of $D$ (Proposition~\ref{prop:ellbasic}(a)). We show $\Delta=0$.

\emph{Small and large.} $\Delta(0)=1-g-0-1+g=0$, since $\ell(W)=g$. For $d>n-3$, $\ell(W-E_d)=0$ (negative degree) and $\ell(E_d)=dn+1-g$
(Proposition~\ref{prop:countforms}), so $\Delta(E_d)=0$.

\emph{Monotonicity.} Let $D\ge0$ and $P$ a point. By Proposition~\ref{prop:ellbasic}(b), $\ell(D+P)-\ell(D)$ and $\ell(W-D)-\ell(W-D-P)$ are $0$ or $1$, and by
Lemma~\ref{lem:reduction} they are not both $1$. Since
\[
\Delta(D+P)-\Delta(D)=\bigl[\ell(D+P)-\ell(D)\bigr]+\bigl[\ell(W-D)-\ell(W-D-P)\bigr]-1,
\]
$\Delta(D+P)\le\Delta(D)$. Adding points one at a time, $\Delta(D_2)\le\Delta(D_1)$ whenever $D_2\ge D_1\ge0$.

\emph{Effective divisors.} Let $D\ge0$. As in the proof of Proposition~\ref{prop:riemann}, there is a form $G$ of some degree $d>n-3$ with
$\operatorname{div}(G)=D+D_0\ge D$, and $\operatorname{div}(G)\equiv E_d$. So $0=\Delta(0)\ge\Delta(D)\ge\Delta(D+D_0)=\Delta(E_d)=0$.

\emph{All divisors.} If $\ell(D)>0$, then $D\equiv D+\operatorname{div}(z)\ge0$ for some $z$, and $\Delta(D)=0$. Note that
$\Delta(W-D)=\ell(W-D)-\ell(D)-(2g-2-\deg D)-1+g=-\Delta(D)$; so if $\ell(W-D)>0$, $\Delta(D)=-\Delta(W-D)=0$. If $\ell(D)=\ell(W-D)=0$, Riemann's inequality
for $D$ and for $W-D$ gives $\deg D\le g-1$ and $2g-2-\deg D\le g-1$, so $\deg D=g-1$ and $\Delta(D)=0-0-(g-1)-1+g=0$.
\end{proof}

\begin{beauty}
The left side of the Riemann--Roch formula is about functions, an analytic or algebraic notion; the right side is about a degree and a
topological invariant. And the correction term $\ell(W-D)$ is itself a space of functions, for the ``dual'' divisor $W-D$, so the formula is a
symmetry: replacing $D$ by $W-D$ changes the sign of both sides. In the proof, the analytic content of Riemann's original argument (residues of
differentials) is replaced by B\'ezout's theorem on a line: a curve of degree $n-2$ cannot meet a line $n-1$ times.
\end{beauty}

Riemann--Roch says most when $W-D$ has negative degree: if $\deg D>2g-2$, then $\ell(D)=\deg D+1-g$ exactly. For genus $1$, $W=0$, and the
consequences are the foundation of Chapter~15.

\begin{corollary}[curves of genus 1]\label{cor:genus1}
Let $C$ be a smooth plane cubic.
\begin{enumerate}[label=(\alph*)]
\item $\ell(D)=\deg D$ for every divisor of degree $\ge1$.
\item If $P\ne Q$, then $P\not\equiv Q$: no function on $C$ has a single simple pole.
\item Fix $O\in C$. Every divisor of degree $0$ is linearly equivalent to $P-O$ for exactly one point $P$.
\item (Nine-point theorem.) Let $G_1,G_2$ be cubics with $\operatorname{div}(G_1)=P_1+\dots+P_8+P_9$ and $\operatorname{div}(G_2)=P_1+\dots+P_8+Q$ (points
counted with multiplicity). Then $Q=P_9$.
\end{enumerate}
\end{corollary}
\begin{proof}
(a) $g=1$, $W=E_0=0$, and $\ell(-D)=0$ when $\deg D\ge1$. (b) If $Q-P=\operatorname{div}(z)$ then $z\in L(P)$, which consists of the constants by (a):
$\ell(P)=1$. Then $\operatorname{div}(z)=0$ and $Q=P$. (c) $\ell(D+O)=1$, so there is a
$z\ne0$ with $\operatorname{div}(z)+D+O=P\ge0$ of degree $1$, a single point; and if $P-O\equiv P'-O$, then $P=P'$ by (b). (d) $Q-P_9=\operatorname{div}(G_2/G_1)$, so
$Q=P_9$ by (b).
\end{proof}

Part (c) identifies the points of $C$ with the group of divisor classes of degree $0$, and this is the group law of Chapter~15; part (d) is what
makes it associative there. Part (b) is the precise sense in which the smooth cubic of Example~\ref{exa:threecubics} is not rational: on the nodal cubic, the slope $t$ of the lines
through the node is a function taking each value once, with a single simple pole, and nothing like it exists on a smooth cubic.

\begin{remark}
Riemann--Roch holds for every smooth projective curve, not only plane ones, with $g$ defined by the formula itself (for instance as $\ell(W)$). Most curves
are not smooth plane curves: a smooth plane curve has genus $0,1,3,6,10,\dots$, never $2$. The curve $y^2=f(x)$, with $f$ of degree $5$ or $6$ without
repeated roots, has genus $2$; its plane model is singular at infinity. For a plane curve with ordinary singularities of multiplicities $r_P$, Fulton's
book proves Riemann--Roch for its smooth model by the same method, with \emph{adjoint} curves (curves with multiplicity $\ge r_P-1$ at each singular point),
and finds $g=\frac{(n-1)(n-2)}2-\sum_P\frac{r_P(r_P-1)}2$. A nodal cubic has genus $0$, consistently with its rational parametrisation.
\end{remark}

\section{Codes from curves}\label{sec:goppa}

A \emph{linear code} of length $N$ over $\mathbb F_q$ is a subspace $\mathcal C\subseteq\mathbb F_q^N$, of dimension $\kappa$ and minimum distance
$\delta=\min\{\text{number of non-zero coordinates of }c:c\in\mathcal C\setminus0\}$. A message of $\kappa$ symbols is encoded as a codeword of $N$ symbols;
any $\delta-1$ erased symbols, or any $\lfloor(\delta-1)/2\rfloor$ wrong ones, can be recovered. One wants $\kappa$ and $\delta$ both large, but the
\emph{Singleton bound} $\kappa+\delta\le N+1$ forbids too much (project a codeword onto its first $\kappa-1$ coordinates: two codewords agree there, and their
difference has weight $\le N-\kappa+1$). Reed--Solomon codes reach it: evaluate the polynomials of degree $<\kappa$ at $N$ distinct points of $\mathbb F_q$;
a non-zero one has at most $\kappa-1$ roots, so $\delta\ge N-\kappa+1$. The catch is $N\le q$: there are only $q$ points on a line.

V.~D.~Goppa's idea (1981) was to evaluate functions on a curve instead of a line, since a curve over $\mathbb F_q$ can have many more than $q$ points. Here is
the version that needs nothing beyond B\'ezout. Let $F$ be a form over $\mathbb F_q$ defining a smooth curve $C$ of degree $n$ and genus $g$ over
$\overline{\mathbb F}_q$, and let $P_1,\dots,P_N$ be points of $C$ with coordinates $(x_i:y_i:1)$ in $\mathbb F_q$. For $d\ge0$ let
\[
\mathcal C_d=\bigl\{\bigl(G(x_1,y_1,1),\dots,G(x_N,y_N,1)\bigr):G\text{ a form of degree }d\text{ over }\mathbb F_q\bigr\}.
\]

\begin{proposition}[Goppa codes from plane curves]\label{prop:goppa}
If $dn<N$, then $\mathcal C_d$ has dimension $\kappa=\dim V_d$ (over $\mathbb F_q$), which is $dn+1-g$ for $d\ge n-2$, and minimum distance
$\delta\ge N-dn$. In particular $\kappa+\delta\ge N+1-g$.
\end{proposition}
\begin{proof}
If $F\nmid G$, then $G$ and $F$ have no common component, and $G$ vanishes at no more than $\sum_PI(P,F\cap G)=dn$ points of $C$ (Theorem~\ref{thm:bezout}
over $\overline{\mathbb F}_q$). So the codeword of $G$ has at least $N-dn>0$ non-zero coordinates; and if $F\mid G$, it is $0$. Hence the kernel of
$G\mapsto$ codeword is exactly $F\cdot R_{d-n}$ (divisibility by $F$ over $\overline{\mathbb F}_q$ and over $\mathbb F_q$ are the same, since dividing is
solving a linear system with coefficients in $\mathbb F_q$), and the dimension is $\dim V_d$, computed over $\mathbb F_q$ exactly as in
Proposition~\ref{prop:countforms}.
\end{proof}

The genus is the price paid with respect to Singleton. For the Hermitian curve $y^4z+yz^4=x^5$ over $\mathbb F_{16}$ (Exercise~\ref{ex:hermitian}) one
gets $N=64$ affine points, $g=6$, and codes of length $64$ over an alphabet of $16$ letters, four times longer than any Reed--Solomon code over
$\mathbb F_{16}$; they can be constructed explicitly and used to correct erasures. Better still, Tsfasman, Vl\u{a}du\c{t} and Zink showed in 1982, using modular curves
(Chapter~17) with about $(\sqrt q-1)g$ points over $\mathbb F_q$, that for $q\ge49$ a square, codes from curves beat the Gilbert--Varshamov bound, which
coding theorists had believed optimal for thirty years. How many points a curve over $\mathbb F_q$ can have is the subject of Chapter~16.

\section{Exercises}

\begin{exercise}[$\star$]\label{ex:circles}
Show that every circle $(x-a)^2+(y-b)^2=r^2$ passes through the two \emph{circular points} $(1:i:0)$ and $(1:-i:0)$ of $\mathbb P^2(\C)$. Deduce that two
distinct circles meet in at most two points of the affine plane, and find all four B\'ezout points of $x^2+y^2=1$ and $(x-1)^2+y^2=1$.
\end{exercise}

\begin{exercise}[$\star$]\label{ex:singcubic}
Let $F$ be an irreducible cubic with a singular point $P$. Show that $P$ is the only singular point and that every line through $P$, other than the
tangents at $P$, meets $F$ in exactly one more point. Deduce that $F$ is rational, and find the parametrisation of $y^2=x^3+x^2$ by lines through the
origin.
\end{exercise}

\begin{exercise}[$\star$]\label{ex:pascal}
(Pascal, 1640.) Let $A_1,\dots,A_6$ be six distinct points on an irreducible conic $Q$, and let $L_i$ be the line $A_iA_{i+1}$ (indices mod $6$). Assume
that the three points $R_1=L_1\cap L_4$, $R_2=L_2\cap L_5$, $R_3=L_3\cap L_6$ are defined (no two opposite sides parallel or equal) and do not lie on $Q$.
Prove that $R_1,R_2,R_3$ are collinear. (Consider the cubics $L_1L_3L_5+\lambda L_2L_4L_6$.)
\end{exercise}

\begin{exercise}[$\star\star$]\label{ex:quartic}
Let $C$ be a smooth plane quartic. Show that $g=3$, that the canonical divisors $W'\equiv W$ with $W'\ge0$ are exactly the divisors $\operatorname{div}(M)$ of the
lines $M$, and that $\ell(P+Q)=1$ for all points $P,Q$: there is no function on $C$ with only two simple poles (a smooth quartic is not
\emph{hyperelliptic}).
\end{exercise}

\begin{exercise}[$\star\star$]\label{ex:weierstrass}
Let $\operatorname{char}k\ne2,3$ and let $C$ be the smooth cubic $Y^2Z=X^3+aXZ^2+bZ^3$, $O=(0:1:0)$, $x=X/Z$, $y=Y/Z$. Show that $\operatorname{div}(Z)=3O$, that
$x$ and $y$ have a pole of order $2$ and $3$ at $O$ and no other pole, and that for $m\ge1$ the functions $x^iy^j$ with $j\in\{0,1\}$, $2i+3j\le m$ form a
basis of $L(mO)$. Explain why seven functions in $L(6O)$ must satisfy a linear relation, and identify it.
\end{exercise}

\begin{exercise}[$\star\star$]\label{ex:conic}
Let $C$ be a smooth conic and $P\in C$. Show that $\ell(P)=2$, that a non-constant $z\in L(P)$ takes every value of $k\cup\{\infty\}$ exactly once on $C$, and that
it can be taken to be the slope of the line through $P$. For $x^2+y^2=1$ and $P=(-1,0)$, write the inverse map and deduce the formula for all primitive
Pythagorean triples.
\end{exercise}

\begin{exercise}[$\star\star\star$]\label{ex:hermitian}
Let $q$ be a power of a prime $p$ and let $H$ be the Hermitian curve $Y^qZ+YZ^q=X^{q+1}$ over $\mathbb F_{q^2}$.
\begin{enumerate}[label=(\alph*)]
\item Show that $H$ is smooth over $\overline{\mathbb F}_p$, of genus $\frac{q(q-1)}2$, with a single point at infinity.
\item Show that $H$ has exactly $q^3$ affine points with coordinates in $\mathbb F_{q^2}$. (Use $x^{q+1}\in\mathbb F_q$ and count solutions of $y^q+y=c$.)
\item Give the parameters of the code $\mathcal C_d$ of Proposition~\ref{prop:goppa} for $q=4$, and compare the number of $\mathbb F_{q^2}$-points with
$q^2+1+2gq$.
\end{enumerate}
\end{exercise}

\section{Hints}

\begin{hint}{ex:circles}
Homogenise and set $Z=0$. For the four points, subtract the two equations.
\end{hint}
\begin{hint}{ex:singcubic}
A line through two singular points meets the cubic with total multiplicity $\ge4$. A line through $P$ meets $F$ with multiplicity $\ge2$ at $P$.
\end{hint}
\begin{hint}{ex:pascal}
Choose $\lambda$ so that the cubic passes through a seventh point of $Q$, then count intersections with $Q$.
\end{hint}
\begin{hint}{ex:quartic}
For $W$, use Lemma~\ref{lem:noetherrep}. For $\ell(P+Q)$, apply Riemann--Roch and count the lines through $P$ and $Q$, or tangent at $P$ if $P=Q$.
\end{hint}
\begin{hint}{ex:weierstrass}
At $O$, use the chart $Y=1$. For the basis, compare dimensions with Corollary~\ref{cor:genus1}(a) and look at pole orders.
\end{hint}
\begin{hint}{ex:conic}
Use Riemann--Roch with $g=0$. The function $z$ defines a map to $\mathbb P^1$; how many points has $\operatorname{div}(z-a)$?
\end{hint}
\begin{hint}{ex:hermitian}
Compute the three partial derivatives in characteristic $p$. For (b), the norm map $x\mapsto x^{q+1}$ goes from $\mathbb F_{q^2}$ to $\mathbb F_q$, and $y\mapsto
y^q+y$ is $\mathbb F_q$-linear from $\mathbb F_{q^2}$ to $\mathbb F_q$.
\end{hint}

\section{Solutions}

\begin{solution}{ex:circles}
\idea Points at infinity of a curve are the zeros of its top-degree part.
\writeup Homogenised, the circle is $(X-aZ)^2+(Y-bZ)^2=r^2Z^2$; on $Z=0$ it reduces to $X^2+Y^2=0$, that is $(X+iY)(X-iY)=0$: the points $(1:i:0)$ and
$(1:-i:0)$, whatever $a,b,r$. Two distinct circles have no common component (a circle is an irreducible conic: $x^2+y^2-r^2$, after translation, has no
linear factor when $r\ne0$), so they meet in $4$ points with multiplicity, two of which are the circular points: at most $2$ affine points remain. For
$x^2+y^2=1$ and $(x-1)^2+y^2=1$, subtracting gives $2x-1=0$, so $x=\frac12$, $y=\pm\frac{\sqrt3}2$; with $(1:\pm i:0)$, these are the four points, each of
multiplicity $1$ (four distinct points and total $4$).
\end{solution}

\begin{solution}{ex:singcubic}
\idea Bézout with a line is a degree count, and singular points eat two intersections each.
\writeup If $L$ is a line through $P$ and $L\nmid F$, then $I(P,F\cap L)\ge2$: in coordinates with $P=(0,0)$ and $L=\{y=0\}$, $F(x,0)$ has no constant or
linear term, as $m_P(F)\ge2$, so $(F,y)=(F(x,0),y)\subseteq(x^2,y)$ and the quotient has dimension $\ge2$. If $P'$ were a second singular point, the line
$PP'$ would meet $F$ with total multiplicity $\ge4>3$, so it would be a component of $F$, contradicting irreducibility. Now let $L$ be a line through $P$,
not tangent at $P$. Then $I(P,F\cap L)=2$ exactly: the lowest-degree part $F_2$ of $F$ at $P$ (of degree $2$; it is not $3$, else $F$ would be a product of
three lines through $P$) does not vanish on the direction of $L$, so $F(x,0)=cx^2+\dots$ with $c\ne0$. By B\'ezout, $L$ meets $F$ in exactly one more point
$\varphi(L)$, with multiplicity $1$. Conversely each point $Q\ne P$ of $F$ lies on the line $PQ$. So the lines through $P$, a family $y=tx$
parametrised by $t\in k$ (plus $x=0$), parametrise $F$, up to finitely many exceptions. For $y^2=x^3+x^2$ and $y=tx$: $t^2x^2=x^3+x^2$, so $x=t^2-1$
(the double root $0$ is $P$), and $(x,y)=(t^2-1,t(t^2-1))$. The tangents at $P$ are $y=\pm x$, the values $t=\pm1$, both giving $P$ itself.
\end{solution}

\begin{solution}{ex:pascal}
\idea A cubic meeting a conic in seven points must contain it, and what is left is a line.
\writeup Let $C_1=L_1L_3L_5$ and $C_2=L_2L_4L_6$. Each $A_i$ lies on one odd-numbered line and one even-numbered line ($A_i\in L_{i-1}\cap L_i$), so both cubics
pass through all six $A_i$; both also pass through $R_1,R_2,R_3$. Take a point $A_7\in Q$ different from the $A_i$, and not on $C_2$ (only finitely many
points of $Q$ are on $C_2$, by B\'ezout, as $Q$ is not a component of $C_2$). Choose $\lambda$ with $C=C_1+\lambda C_2$ vanishing at $A_7$; this is possible
since $C_2(A_7)\ne0$. Then $C$ is a cubic (non-zero: $C_1$ and $C_2$ are not proportional, having different line factors) vanishing at seven distinct points of
the conic $Q$. If $Q$ were not a component of $C$, B\'ezout would give at most $6$ intersection points; so $C=QM$ for a line $M$. Now $C$ vanishes at each $R_j$,
and $R_j\notin Q$, so $M(R_j)=0$: the three points lie on the line $M$.
\end{solution}

\begin{solution}{ex:quartic}
\idea For a quartic, $W=E$: canonical divisors are line sections, and two points determine a line.
\writeup $n=4$, so $g=\frac{3\cdot2}2=3$ and $W=E_1$. By Lemma~\ref{lem:noetherrep} with $D=0$, $L(E_1)=\{M/L:M\text{ a linear form}\}$, and the effective
divisors linearly equivalent to $W$ are the $\operatorname{div}(z)+E_1=\operatorname{div}(M)$ for $z=M/L\ne0$: the line sections. For $\ell(P+Q)$: Riemann--Roch gives
$\ell(P+Q)=2+1-3+\ell(W-P-Q)=\ell(E_1-(P+Q))$, and by Lemma~\ref{lem:noetherrep} this is the dimension of the space of linear forms $M$ with
$\operatorname{div}(M)\ge P+Q$. If $P\ne Q$, the condition is $M(P)=M(Q)=0$, and the linear forms vanishing at two distinct points form a space of dimension $1$.
If $P=Q$, the condition is $I(P,C\cap M)\ge2$, which says $M$ is the tangent line at $P$ (a line through $P$ has $I=1$ iff it is not the tangent, by
Proposition~\ref{prop:smoothdvr}: a non-tangent line is a uniformiser); again dimension $1$. So $\ell(P+Q)=1$: $L(P+Q)$ is the constants.
\end{solution}

\begin{solution}{ex:weierstrass}
\idea Pole orders $2i+3j$ are all distinct when $j\le1$, which makes the functions independent; Riemann--Roch says there are exactly enough.
\writeup Setting $Z=0$ gives $X^3=0$: the only point at infinity is $O$, and $\deg\operatorname{div}(Z)=3$, so $\operatorname{div}(Z)=3O$. In the chart $Y=1$, with
$u=X/Y$, $v=Z/Y$, the curve is $v=u^3+auv^2+bv^3$, smooth at $(0,0)$ with tangent $v=0$; $u$ is a uniformiser, and $v=u^3\cdot(\text{unit})$, so
$\operatorname{ord}_O(v)=3$. Then $x=X/Z=u/v$ has order $1-3=-2$ and $y=Y/Z=1/v$ has order $-3$ at $O$. Away from $O$, $Z\ne0$ and $x,y$ are regular. The
functions $x^iy^j$ ($j\le1$) have pole order exactly $2i+3j$ at $O$, and these orders are distinct, so the functions are linearly independent (in a
relation, the term with the largest pole order cannot cancel). Those with $2i+3j\le m$ lie in $L(mO)$, and there are exactly $m$ of them for
$m\ge1$: the values $0,2,3,4,\dots,m$ of $2i+3j$, all integers in $[0,m]$ except $1$. Since $\ell(mO)=m$ (Corollary~\ref{cor:genus1}(a)), they form a basis.
In $L(6O)$ the seven functions $1,x,y,x^2,xy,x^3,y^2$ are dependent, as $\ell(6O)=6$; the relation is the equation of the curve, $y^2=x^3+ax+b$, in which
$y^2$ and $x^3$ are the two functions of pole order $6$, and must cancel.
\end{solution}

\begin{solution}{ex:conic}
\idea Genus $0$ makes $\ell(P)=2$; a function with one simple pole is a bijection to the projective line.
\writeup $n=2$, $g=0$, and $\deg(W-P)=-3<0$, so $\ell(P)=1+1-0=2$. Take $z\in L(P)$ non-constant; it has a simple pole at $P$ (a pole of order $\ge1$
bounded by $1$), so $\deg\operatorname{div}(z)=0$ forces $\operatorname{div}(z)=Q-P$ for one point $Q$. Likewise for each $a\in k$, $z-a\in L(P)$, so
$\operatorname{div}(z-a)=Q_a-P$: the value $a$ is taken exactly once, at $Q_a$, and $\infty$ only at $P$. With $M$ a line through $P$, not tangent, and $T$
the tangent at $P$, $z=M/T$ works: $\operatorname{div}(M)-\operatorname{div}(T)=(P+Q)-2P$. For the circle and $P=(-1,0)$, $T=x+1$, $M=y$: $z=y/(x+1)=t$, the
slope. Then $y=t(x+1)$ and $x^2+t^2(x+1)^2=1$, that is $(x+1)\bigl((1+t^2)x-(1-t^2)\bigr)=0$, so $x=\frac{1-t^2}{1+t^2}$, $y=\frac{2t}{1+t^2}$. Since $z$ is a
bijection, $t\in\Q$ gives exactly the rational points other than $P$. With $t=\frac uv$ in lowest terms, $(x,y)=\bigl(\frac{v^2-u^2}{v^2+u^2},\frac{2uv}{v^2+u^2}\bigr)$,
and clearing denominators, the primitive solutions of $a^2+b^2=c^2$ are, up to order and signs, $(v^2-u^2,2uv,v^2+u^2)$ with $\gcd(u,v)=1$ and $u\not\equiv v\pmod2$
(when $u,v$ are both odd, everything is divisible by $2$, and dividing gives the same form with the roles of $a$ and $b$ swapped).
\end{solution}

\begin{solution}{ex:hermitian}
\idea In characteristic $p$, most terms of the partial derivatives die, which makes smoothness easy; the point count is about the norm and trace from
$\mathbb F_{q^2}$ to $\mathbb F_q$.
\writeup (a) With $F=Y^qZ+YZ^q-X^{q+1}$: $F_X=-(q+1)X^q=-X^q$, $F_Y=qY^{q-1}Z+Z^q=Z^q$, $F_Z=Y^q+qYZ^{q-1}=Y^q$, since $q=0$ in $k$. They vanish together only at
$X=Y=Z=0$, so $H$ is smooth (Lemma~\ref{lem:projsmooth}), of degree $n=q+1$ and genus $\frac{q(q-1)}2$. On $Z=0$, $X^{q+1}=0$: only $(0:1:0)$.

(b) For $x\in\mathbb F_{q^2}$, $c=x^{q+1}$ satisfies $c^q=x^{q^2+q}=x^{1+q}=c$ (as $x^{q^2}=x$), so $c\in\mathbb F_q$. The map $T(y)=y^q+y$ is $\mathbb F_q$-linear
from $\mathbb F_{q^2}$ to $\mathbb F_q$ (by the same computation, $T(y)^q=y^{q^2}+y^q=T(y)$), and its kernel, the roots of $y^q+y$, has at most $q$ elements, so its image
has at least $q^2/q=q$ elements: $T$ is onto $\mathbb F_q$, with fibres of size exactly $q$. Hence for each of the $q^2$ values of $x$ there are exactly $q$
values of $y$: $q^3$ affine points.

(c) For $q=4$: $n=5$, $g=6$, $N=64$ over $\mathbb F_{16}$. For $3\le d\le12$ ($5d<64$), $\mathcal C_d$ has dimension $5d-5$ and minimum distance $\ge64-5d$; for
example $d=6$ gives a $[64,25,\ge34]$ code, and $d=10$ a $[64,45,\ge14]$ code. The curve has $q^3+1=65$ points over $\mathbb F_{16}$, and
$q^2+1+2gq=16+1+48=65$: the Hermitian curve attains the bound of Chapter~16 exactly (it is \emph{maximal}).
\end{solution}

\section*{Sources and further reading}

\begin{itemize}
\item W.~Fulton, \emph{Algebraic Curves: An Introduction to Algebraic Geometry}, 2008 edition, freely available from the author's web page. Chapters~2--5 and~8
are the route followed here (intersection numbers, B\'ezout, Noether, Riemann--Roch), with the general case of plane curves with ordinary singularities.
\item E.~Brieskorn and H.~Kn\"orrer, \emph{Plane Algebraic Curves}, Birkh\"auser, 1986; reprinted 2012. History, pictures, and the geometry of real curves.
\item I.~R.~Shafarevich, \emph{Basic Algebraic Geometry 1}, 3rd ed., Springer, 2013, Chapters~III--IV: Riemann--Roch for all smooth curves.
\item D.~Eisenbud, M.~Green and J.~Harris, ``Cayley--Bacharach theorems and conjectures'', \emph{Bull.\ AMS} 33 (1996), 295--324. From Pappus and Pascal to modern
algebra, starting with the nine-point theorem.
\item H.~Stichtenoth, \emph{Algebraic Function Fields and Codes}, 2nd ed., Springer GTM~254, 2009. Riemann--Roch for function fields over any field, Goppa
codes, the Hermitian curve.
\item M.~A.~Tsfasman, S.~G.~Vl\u{a}du\c{t} and T.~Zink, ``Modular curves, Shimura curves, and Goppa codes, better than Varshamov--Gilbert bound'',
\emph{Math.\ Nachr.} 109 (1982), 21--28.
\item \emph{History.} Maclaurin (1720), Cramer (1750) and B\'ezout (1779): Brieskorn--Kn\"orrer, above, Chapter~1, and
the Eisenbud--Green--Harris article. Riemann (1857), Roch (1865) and Brill--Noether (1874): J.~J.~Gray, ``The
Riemann--Roch theorem and geometry, 1854--1914'', \emph{Doc.\ Math.}, Extra Volume ICM~III (1998), 811--822; Roch's
death at twenty-six: J.~J.~O'Connor and E.~F.~Robertson, ``Gustav Roch'', MacTutor History of Mathematics archive.
\end{itemize}
